As a final validation, the Kalman filter model with the neural network is tested.

Evaluating the neural network in the ensemble Kalman filter method presents a more intricate task.
Due to the computationally involved ETKF algorithm, a quick computation of the error metrics for a larger dataset is not feasible.
As the main focus of this work was on integrating the neural network as the state model operator, a simple observation model with simulated missing values was used.
However, as already seen with the COT dataset, missing valued present a significant problem in many real data sources.
Hence, the simulated exercise is not unrealistic, even though the model could be capable to more sophisticated inverse problems such as data fusion. 

Similar COT images outside of the training dataset of $128 \times 128$ pixels and similar synthetic examples are used as with evaluations of the deterministic model.
Despite using images of the same shape as the neural network works with, the model could work with different shape of images.
Larger images could be downsampled for the neural network model to create a state operator to work with larger state vectors in the ETKF or vice versa.
In order to visualise the performance of the motion models, the presented example images are selected such that they largely adhere to the advection equation with some random variation to demonstrate observation assimilation.

For fine-tuning the ETKF model presented in Subsection~\ref{subsec:etkfmodel}, a few assumptions were made related to the objectives of the verification.
The ETKF parameters were chosen to give a small uncertainty for the observations.
The intention was to return the observation values where available and, otherwise, values warped from the previous state.
This choice stems from the quickly evolving COT data not only within the model's capability but also to verify the results better.
In real use of the model, these parameters are subject to case-by-case tuning depending on the used data.
Particularly, the observation uncertainty should be a realistic uncertainty estimate of the measurements instead of the assumed constant of 0.23 used for these examples.
The data was scaled to $[0, 1]$ range similarly as when training the network.
For the initial ensemble and inflation, spatially varying additive noise was generated over the entire image area and then multiplied pixelwise with the image at hand.
Localisation was performed in $16 \times 16$ pixel blocks.
An ensemble size of 500 is used as it appears to give a good balance between results and computing time.

Simulated synthetic data provides a quick visual validation of the ETKF model.
The simulated observed noisy data in Figure~\ref{fig:etkfexamplesynthetic} has a noisy background with rectangular blocks moving according to a uniform velocity field.
A vertical block of observations has been removed from the middle of the observation images.
In this example, the ETKF model works as expected by filling in the missing values and removing fine noise from the observations.
The neural network model appears to overestimate the velocity fields a little, but the kalman filter corrects the estimates where observations are available.
As a new rectangle appears in the observations, it also appears in the filtered images as expected.
This example thus suggests that the model works as intended.

\begin{figure}[H]
    \inputpypgf{experiments/fig/etkf_blocks.pypgf}
    \caption{\label{fig:etkfexamplesynthetic}ETKF model example with synthetic data.
    The second row represents the observed data with a vertical black block of missing values, and the first row has the ensemble means of the estimated filtering distribution with the predicted velocity fields.}
\end{figure}

For verification with the COT data, a sequence of images was selected to compare the neural network model as the state operator against the Lucas-Kanade method and the identity operator.
In Figure~\ref{fig:etkfexamplecot}, the filtered images of the three different state operators are presented.
As expected, due to the small observation error assumption, the used state propagator shows relatively small differences in the areas of the image where observations are available.
However, in the area of missing observations, the motion models provide a clear advantage, as the identity operator leaves the values to the initial state throughout the time steps.

\begin{figure}[h]
    \inputpypgf{experiments/fig/etkf_cot.pypgf}
    \caption{\label{fig:etkfexamplecot}A comparison of different model operators for the ETKF\@.
        The leftmost column contains the observations with a vertical black block of simulated missing observations; the second column contains the filtering distribution with the neural network as the model operator, the third column with the Lucas-Kanade method, and the third with a trivial identity operator.
    Estimated velocity fields using the corresponding models in the filtered images show similar patterns with the Lucas-Kanade method and the neural network model.}
\end{figure}

As the Kalman filter produces probabilistic estimates, another way to interpret the results is through the uncertainty of the forecast ensembles, observation uncertainty, and posterior ensemble uncertainties.
Figure~\ref{fig:etkferrorbars} shows how the ETKF corrects the mean deviations with the observations for the neural network images of Figure~\ref{fig:etkfexamplecot}.
Figure~\ref{fig:etkfstd}, shows the state ensemble pixelwise standard deviation across the same time steps.
The uncertainty estimates are mostly a result of the selected inflation technique and suitable observation standard deviation parameter.
Particularly, the pixelwise uncertainties in Figure~\ref{fig:etkfstd} show that there is larger uncertainty in the area of clouds which is a direct consequence of the added noise fields.

\begin{figure}[h]
    \inputpypgf{experiments/fig/errorbars.pypgf}
    \caption{\label{fig:etkferrorbars}Mean uncertainty estimates of the neural network computed prior forecast, observations, and posterior estimates given by the ETKF of the example in Figure~\ref{fig:etkfexamplecot}.
        The error bar widths correspond to two times the spatial mean of the standard deviations.
    For the observation, the standard deviation was assumed to be 0.23.}
\end{figure}

\begin{figure}[h]
    \inputpypgf{experiments/fig/etkfstd.pypgf}
    \caption{\label{fig:etkfstd}Two standard deviations of the neural network ETKF model of the images in Figure~\ref{fig:etkfexamplecot}.
    Missing observations lead to higher uncertainty, and also, due to the use of inflation techniques, uncertainty is larger in the areas with more clouds.}
\end{figure}

As seen in these examples, a significant difficulty in the uncertainty quantification with the Kalman filter is that the filter parameters effectively control the magnitude of uncertainty.
Especially with the ensemble methods, where the ensemble essentially defines the state propagator model uncertainty, many unknown parameters must be guessed or estimated by hand.
Particularly, the inflation method plays a crucial role accurate uncertainty estimates as seen with these examples.
In wider use of the model, these parameters could be better estimated with some statistical methods instead of hand-tuning with the objective to obtain visually pleasing results.

The inflation and initial ensemble generation are even more critical in possible forecast usage, as no observations are available to support the forecasts.
In addition, as the computation of a forecast distribution in the ensemble filter is not limited by the standard Kalman filter assumptions, the range of possibilities for ensemble generation and inflation is much broader.
In a recent work related to the forecasting problem presented in~\cite{intraday}, an ensemble is generated by multiplying the Fourier transforms of a noise field and a state estimate.
In that work, an advection model is also combined with an autoregressive model that could potentially improve the predictive performance of the neural network of this work for modeling the COT data.
